\section{Conclusiones}

La realización de esta práctica ha permitido profundizar en el comportamiento y la aplicabilidad de diferentes de diferentes metaheurísticas.

En lo relativo al algoritmo de \textbf{Optimización Basada en Biogeografía (BBO)}, su implementación nos ha resultado notablemente intuitiva. El mayor valorq que hemos aprendido de esta fase ha sido la experimentación con las funciones de prueba $T_1$ a $T_6$; someter al algoritmo a estos entornos controlados nos ha permitido comprender cómo se comporta la búsqueda ante Frentes de Pareto con geometrías muy diversas (convexas, cóncavas o discontinuas), revelando la gran utilidad de estos \textit{benchmarks} estándar para validar la robustez de cualquier optimizador.

Por su parte, el \textbf{Sistema de Hormigas Elitista} ha destacado por la elegancia y originalidad de su inspiración biológica basada en el rastro de feromonas. Inicialmente, la resolución del problema mono-objetivo planteado en el enunciado nos pareció sencilla debido a la eficacia del algoritmo. Es por esto que decidimos explorar los límites de la técnica, extendiendo la implementación hacia una variante multi-objetivo. Esta decisión aumentó la complejidad técnica y nos permitió observar cómo funcionan las soluciones de compromiso en problemas combinatorios.

En cuanto a la metaheurística \textbf{FANS (Fuzzy Adaptive Neighborhood Search)}, lo más reseñable ha sido comprender su gran polivalencia. El uso de lógica difusa para evaluar las soluciones brinda al algoritmo de una flexibilidad única, permitiéndole adaptarse a un gran abanico de problemas de optimización sin necesidad de redefinir rígidamente los criterios de aceptación, lo cual lo convierte en una herramienta sumamente versátil en comparación con métodos de búsqueda local más tradicionales.

Finalmente, desde una perspectiva técnica, el desarrollo de estos algoritmos en Python ha supuesto un aprendizaje muy valioso. Hemos comprobado la existencia de una gran cantidad de repositorios de código abierto y librerías dedicadas a la computación evolutiva. La disponibilidad de estas referencias no solo acelera el proceso de implementación, sino que permite contrastar nuestras soluciones con estándares de la comunidad, facilitando la detección de errores y mejorando de la eficiencia computacional de los algoritmos.